\documentclass[12pt,a4paper]{article}
\usepackage[utf8]{inputenc}
\usepackage[T1]{fontenc}
\usepackage{enumitem}
\usepackage[margin=2.5cm]{geometry}
\usepackage{parskip}
\usepackage{amsmath}

\title{midterm 4 — Final Exam}
\date{}

\begin{document}

\maketitle



\noindent
\textbf{Total points:} 5
\hfill
\textbf{Number of questions:} 5

\vspace{1em}
\hrule
\vspace{1.5em}


\noindent
\textbf{Question 1 (1.00 pts).}~
Which of the following is a characteristic of a proposition?


\begin{enumerate}[label=\textbf{\Alph*.}]

  \item It can be true or false

  \item It can only be true

  \item It can only be false

  \item It cannot be determined as true or false

\end{enumerate}



\vspace{0.5em}


\noindent
\textbf{Question 2 (1.00 pts).}~
Which of the following is equivalent to the expression '.x D y/' in terms of membership?


\begin{enumerate}[label=\textbf{\Alph*.}]

  \item 8z: .z 2 x IFF z 2 y/

  \item 8z: .z 2 x IMPLIES z 2 y/

  \item 8z: .z 2 y IFF z 2 x/

  \item 8z: .z 2 x OR z 2 y/

\end{enumerate}



\vspace{0.5em}


\noindent
\textbf{Question 3 (1.00 pts).}~
If the hypothesis of an implication is false, then the implication as a whole is considered true.


\begin{enumerate}[label=\textbf{\Alph*.}]
  \item True
  \item False
\end{enumerate}



\vspace{0.5em}


\noindent
\textbf{Question 4 (1.00 pts).}~
Every propositional formula is equivalent to both a full disjunctive normal form and a full conjunctive normal form.


\begin{enumerate}[label=\textbf{\Alph*.}]
  \item True
  \item False
\end{enumerate}



\vspace{0.5em}


\noindent
\textbf{Question 5 (1.00 pts).}~
A predicate formula is in prenex form when all its quantifiers appear at the end of the formula.


\begin{enumerate}[label=\textbf{\Alph*.}]
  \item True
  \item False
\end{enumerate}





\end{document}
